\subsection{Riesgos Financieros}


\textbf{Naturaleza y Clasificación de los Riesgos Financieros}

\noindent En el ámbito corporativo, el riesgo financiero se entiende como la probabilidad de que una organización enfrente impactos económicos adversos derivados de la volatilidad del mercado, decisiones internas o fenómenos exógenos. Para un análisis adecuado, estos riesgos se categorizan en cuatro grandes grupos:

\vspace{0.3cm}
\begin{itemize}[label=\textcolor{azulp}{\textbullet}, leftmargin=0.8cm]
    \item \textbf{Riesgo de Mercado:} Se origina por movimientos desfavorables en los precios de los activos financieros, englobando la volatilidad en las tasas de interés y los tipos de cambio. Su medición se apoya fuertemente en métricas cuantitativas como el Valor en Riesgo (VaR) y el Valor en Riesgo Condicional (CVaR). Un ejemplo clásico es la depreciación súbita del portafolio accionario de un inversionista.
    
    \item \textbf{Riesgo de Contraparte (o Crediticio):} Se materializa cuando un ente involucrado en un contrato financiero falla en el cumplimiento de sus obligaciones. Este riesgo, típico en el impago de bonos o préstamos, se cuantifica a través de métricas como la Pérdida Esperada, la Pérdida No Esperada y el análisis de la Distribución de Pérdidas.
    
    \item \textbf{Riesgo de Liquidez:} Ocurre cuando una institución carece del capital de trabajo necesario para saldar sus compromisos a corto plazo, viéndose forzada a liquidar activos de manera apresurada y con descuentos significativos respecto a su valor real de mercado.
    
    \item \textbf{Riesgo Operacional:} Deriva de vulnerabilidades o fallas en los sistemas, procesos internos o el capital humano de la empresa, mermando su eficiencia operativa. Este espectro abarca:
    \begin{itemize}[label=\textcolor{azulp}{$\circ$}]
        \item \textit{Riesgo tecnológico:} Fracaso en la obtención de los beneficios económicos esperados tras la implementación de nuevas tecnologías.
        \item \textit{Riesgo país (o soberano):} Obstáculos en los flujos de pago causados por intervenciones o políticas restrictivas de gobiernos extranjeros.
        \item \textit{Riesgo de insolvencia:} Insuficiencia de recursos de capital para amortiguar la depreciación de los activos.
    \end{itemize}
    Para medirlo, se emplean herramientas análogas a las del riesgo crediticio (Pérdida Esperada y No Esperada). Un ejemplo contemporáneo sería la paralización de actividades producto de un ataque cibernético.
\end{itemize}

\vspace{0.5cm}
\noindent \textcolor{azulp}{\textbf{Gestión Integral de Riesgos}}
\vspace{0.2cm}


\noindent Administrar el riesgo no implica eliminarlo, sino identificarlo, medirlo y asegurar que la rentabilidad esperada compense la exposición asumida. Una gestión eficaz optimiza la toma de decisiones y protege el capital de la firma mediante un ciclo continuo de cinco fases:

\begin{enumerate}[label=\textcolor{azulp}{\textbf{\arabic*.}}, leftmargin=0.8cm]
    \item \textbf{Identificación:} Es el diagnóstico inicial de las posibles amenazas. Se nutre de datos históricos, juicio de expertos, checklists y análisis de escenarios. Por ejemplo, detectar la vulnerabilidad de la empresa ante un alza en las tasas de interés.
    
    \item \textbf{Evaluación y Cuantificación:} Consiste en estimar la severidad y la probabilidad de ocurrencia de los riesgos detectados. Las principales metodologías incluyen:
    \begin{itemize}[label=-]
        \item \textit{Valor en Riesgo (VaR):} Proyecta la pérdida máxima bajo un nivel de confianza y horizonte temporal dados.
        \item \textit{Análisis de Sensibilidad:} Evalúa el impacto financiero ante la alteración de variables críticas.
        \item \textit{Simulación Monte Carlo:} Genera trayectorias estocásticas para construir distribuciones empíricas de pérdidas potenciales.
        \item \textit{Modelos Paramétricos:} Utilizan supuestos sobre distribuciones de probabilidad teóricas para estimar resultados.
    \end{itemize}
    
    \item \textbf{Mitigación y Control:} Implementación de tácticas defensivas, tales como:
    \begin{itemize}[label=-]
        \item \textit{Diversificación:} Asignación de capital en múltiples activos no correlacionados.
        \item \textit{Coberturas (Hedging):} Uso de instrumentos derivados (swaps, futuros) para neutralizar volatilidades.
        \item \textit{Transferencia:} Cesión del riesgo a terceros mediante pólizas de seguro.
        \item \textit{Políticas internas:} Fijación de límites estrictos de tolerancia al riesgo.
    \end{itemize}
    
    \item \textbf{Monitoreo Continuo:} Supervisión dinámica del entorno y los portafolios a través de reportes recurrentes, Indicadores Clave de Riesgo (KRI) y revisiones de auditoría.
    
    \item \textbf{Cultura Organizacional:} Consolidación de una mentalidad de prevención y control a nivel institucional. 
\end{enumerate}

\vspace{0.3cm}
\noindent Una administración de riesgos madura no solo minimiza sorpresas negativas, sino que dota a la empresa de mayor resiliencia y estabilidad en la consecución de sus objetivos financieros.
